% ==================================================================
%  SP29-Notes — общая преамбула конспектов
%  Движок: XeLaTeX
%  Подключается из main.tex каждого предмета; не редактируется под предмет
%  (для предметных настроек есть meta.tex).
% ==================================================================

% ---------- язык и шрифты -----------------------------------------
\usepackage{polyglossia}
\setdefaultlanguage{russian}
\setotherlanguage{english}

\usepackage{unicode-math}
% Шрифты Libertinus загружаются по именам файлов: так они находятся
% через kpathsea в TeX Live (пакет libertinus-otf) и не требуют
% установки в систему и локальной папки fonts/ в репозитории.
\setmainfont{LibertinusSerif-Regular.otf}[
  BoldFont       = LibertinusSerif-Bold.otf,
  ItalicFont     = LibertinusSerif-Italic.otf,
  BoldItalicFont = LibertinusSerif-BoldItalic.otf,
  Ligatures=TeX, Numbers=Proportional]
% Гротеск — Fira Sans (пакет fira в TeX Live). Libertinus Sans не берём:
% в его файле Bold лежат засечковые начертания, кириллица уезжает в сериф.
\setsansfont{FiraSans-Regular.otf}[
  BoldFont       = FiraSans-SemiBold.otf,
  ItalicFont     = FiraSans-Italic.otf,
  BoldItalicFont = FiraSans-SemiBoldItalic.otf,
  Ligatures=TeX, Scale=MatchLowercase]
\setmonofont{FiraMono-Regular.otf}[Scale=MatchLowercase]
\setmathfont{LibertinusMath-Regular.otf}

\usepackage{microtype}

% ---------- геометрия ---------------------------------------------
\usepackage[
  a4paper,
  top=24mm, bottom=22mm, left=26mm, right=32mm,
  marginparwidth=20mm, marginparsep=6mm,
  headsep=7mm, footskip=12mm
]{geometry}

\usepackage[skip=6pt plus 2pt, indent=0pt]{parskip}
\linespread{1.04}

% ---------- палитра -----------------------------------------------
\usepackage[dvipsnames]{xcolor}

\definecolor{ink}      {HTML}{1B1B23}   % основной текст
\definecolor{inkSoft}  {HTML}{5C5C6B}   % вторичный текст
\definecolor{ruleC}    {HTML}{DCDCE4}   % линейки
\definecolor{accent}   {HTML}{3B4A8C}   % акцент: номера, ссылки

% Приглушённая палитра, сгруппированная по смыслу:
% утверждения — синяя семья (различаются светлотой), определения — бирюза,
% примеры — зелень, замечания — тёплый серый, вставки составителя — охра.
\definecolor{cThm}     {HTML}{4C5C96}   % теорема
\definecolor{cLem}     {HTML}{6B7BA6}   % лемма
\definecolor{cCor}     {HTML}{7E6FA6}   % следствие
\definecolor{cProp}    {HTML}{5D7A93}   % предложение
\definecolor{cStmt}    {HTML}{5D7A93}   % утверждение
\definecolor{cDef}     {HTML}{3E7C78}   % определение
\definecolor{cEx}      {HTML}{628A6B}   % пример
\definecolor{cRem}     {HTML}{7A7A82}   % замечание
\definecolor{cAdd}     {HTML}{A08050}   % дополнение составителя

\color{ink}

% ---------- математика --------------------------------------------
\usepackage{amsmath, amsthm, mathtools}   % amssymb не нужен: символы даёт unicode-math
\usepackage{bm}

\newcommand{\NN}{\mathbb{N}}
\newcommand{\ZZ}{\mathbb{Z}}
\newcommand{\QQ}{\mathbb{Q}}
\newcommand{\RR}{\mathbb{R}}
\newcommand{\CC}{\mathbb{C}}
\newcommand{\eps}{\varepsilon}
\newcommand{\vphi}{\varphi}
\newcommand{\bigO}{\mathcal{O}}
\newcommand{\set}[1]{\left\{#1\right\}}
\newcommand{\abs}[1]{\left|#1\right|}
\newcommand{\norm}[1]{\left\lVert#1\right\rVert}
\DeclareMathOperator*{\argmin}{arg\,min}
\DeclareMathOperator*{\argmax}{arg\,max}

\setlength{\abovedisplayskip}{7pt}
\setlength{\belowdisplayskip}{7pt}
\setlength{\abovedisplayshortskip}{4pt}
\setlength{\belowdisplayshortskip}{4pt}

% ---------- списки, таблицы, рисунки ------------------------------
\usepackage{enumitem}
\setlist{itemsep=2pt, topsep=4pt, parsep=0pt}
\setlist[enumerate]{label=\arabic*., leftmargin=1.5em}
\setlist[itemize]{label=\textcolor{inkSoft}{\textbullet}, leftmargin=1.3em}

\usepackage{booktabs}
\usepackage{graphicx}
\usepackage{float}
\usepackage{caption}
\usepackage{subcaption}
\captionsetup{
  font={small}, labelfont={sf,bf,color=ink}, textfont={color=inkSoft},
  labelsep=period, skip=6pt
}

\usepackage{tikz}
\usepackage{pgfplots}
\pgfplotsset{compat=1.18}
\usetikzlibrary{arrows.meta, positioning, calc, decorations.pathreplacing,
  fit, backgrounds, patterns, shapes.geometric, matrix, intersections, quotes}
\usepackage{tikz-cd}

% polyglossia делает " активным символом (русские шорткаты вроде "-).
% Это ломает синтаксис подписей tikz-cd: \arrow[r, "f"] превращается
% в неизвестный ключ /tikz/"f". Шорткаты нам не нужны — возвращаем катод.
\AtBeginDocument{\catcode`\"=12\relax}

% ---------- заголовки ---------------------------------------------
\usepackage{titlesec}

\titleformat{\section}[hang]
  {\sffamily\bfseries\Large\color{ink}}
  {\color{accent}\thesection}{0.55em}{}
  [\vspace{-3pt}{\color{ruleC}\titlerule[0.6pt]}]
\titlespacing*{\section}{0pt}{18pt plus 4pt}{9pt}

\titleformat{\subsection}[hang]
  {\sffamily\bfseries\large\color{ink}}
  {\color{accent}\thesubsection}{0.5em}{}
\titlespacing*{\subsection}{0pt}{13pt plus 3pt}{5pt}

\titleformat{\subsubsection}[runin]
  {\sffamily\bfseries\color{ink}}{}{0pt}{}[.\quad]
\titlespacing*{\subsubsection}{0pt}{9pt}{0.6em}

% ---------- колонтитулы -------------------------------------------
\usepackage{fancyhdr}
\pagestyle{fancy}
\fancyhf{}
\renewcommand{\headrulewidth}{0.4pt}
\renewcommand{\headrule}{\hbox to\headwidth{\color{ruleC}\leaders\hrule height \headrulewidth\hfill}}
\fancyhead[L]{\footnotesize\sffamily\color{inkSoft}\subjectshort}
\fancyhead[R]{\footnotesize\sffamily\color{inkSoft}\nouppercase{\leftmark}}
\fancyfoot[C]{\footnotesize\sffamily\color{inkSoft}\thepage}
\fancypagestyle{plain}{\fancyhf{}\renewcommand{\headrulewidth}{0pt}%
  \fancyfoot[C]{\footnotesize\sffamily\color{inkSoft}\thepage}}
\renewcommand{\sectionmark}[1]{\markboth{#1}{}}

% ---------- ссылки ------------------------------------------------
\usepackage{hyperref}
\hypersetup{
  colorlinks=true, linkcolor=accent, urlcolor=accent, citecolor=accent,
  linktoc=all, bookmarksnumbered=true, pdfencoding=unicode
}
\usepackage{bookmark}
\usepackage[nameinlink]{cleveref}

% ==================================================================
%  Окружения
% ==================================================================
\usepackage[most]{tcolorbox}
\tcbuselibrary{breakable, skins, theorems}

% Форма: без рамки, светлая подложка, акцентная планка слева,
% левый край прямой (заподлицо с планкой), правые углы скруглены.
% Заголовок — не «ярлык», а первая строка внутри блока.
\tcbset{
  notebox/.style={
    enhanced, breakable,
    boxrule=0pt, frame hidden,
    sharp corners=west, arc=2mm,
    left=12pt, right=12pt, top=8pt, bottom=8pt,
    before skip=12pt, after skip=12pt,
    fonttitle=\sffamily\bfseries\small,
    detach title,
    separator sign={\ \textperiodcentered\ },
  }
}

% Все утверждения нумеруются одним счётчиком внутри секции:
% Определение 1.1, Теорема 1.2, Лемма 1.3, ... — номер однозначен.
\tcbset{
  boxlook/.style={
    colback=#1!6!white,
    coltitle=#1,
    borderline west={2.4pt}{0pt}{#1},
    before upper={\tcbtitle\par\vspace{3pt}},
  }
}

\newcommand{\declarebox}[4]{%
  % #1 внутреннее имя, #2 подпись, #3 цвет, #4 префикс метки
  \newtcbtheorem[use counter from=Theorem]{#1}{#2}{
    notebox, boxlook=#3,
  }{#4}%
}

\newtcbtheorem[auto counter, number within=section]{Theorem}{Теорема}{
  notebox, boxlook=cThm,
}{thm}
\declarebox{Lemma}      {Лемма}        {cLem} {lem}
\declarebox{Corollary}  {Следствие}    {cCor} {cor}
\declarebox{Proposition}{Предложение}  {cProp}{prop}
\declarebox{Statement}  {Утверждение}  {cStmt}{stmt}
\declarebox{Definition} {Определение}  {cDef} {def}
\declarebox{Example}    {Пример}       {cEx}  {ex}
\declarebox{Remark}     {Замечание}    {cRem} {rem}

% Дополнение составителя: пунктирная рамка, без нумерации
\newtcolorbox{Addition}[1][]{
  notebox, boxlook=cAdd,
  title={\addedglypht\ \,Дополнение\ifstrempty{#1}{}{: #1}},
  fontupper=\small,
}

% Доказательство: без рамки, тонкая линейка слева
\newtcolorbox{prooof}[1][Доказательство]{
  enhanced, breakable, colback=white, colframe=white,
  boxrule=0pt, sharp corners,
  borderline west={0.8pt}{0pt}{ruleC},
  left=9pt, right=0pt, top=3pt, bottom=3pt,
  before skip=9pt, after skip=11pt,
  before upper={\textbf{#1.}\hspace{0.5em}},
}
\newcommand{\qedsym}{\textcolor{inkSoft}{\rule[0.02ex]{1.25ex}{1.25ex}}}
\renewenvironment{proof}[1][Доказательство]
  {\begin{prooof}[#1]}{\hfill\qedsym\end{prooof}}

% ---------- пометка «дописано составителем» -----------------------
\usepackage{marginnote}
\renewcommand*{\marginfont}{\footnotesize\sffamily\color{cAdd}}

\newcommand{\addedglyph}{%
  \tikz[baseline=-0.55ex]{
    \node[draw=cAdd, line width=0.6pt, circle, inner sep=1.1pt,
          font=\tiny\sffamily\bfseries, text=cAdd] {+};}%
}
% \dop        — значок на полях
% \dop[текст] — значок с короткой подписью
\newcommand{\addedglypht}{%
  \tikz[baseline=-0.55ex]{
    \node[draw=cAdd, line width=0.7pt, circle, inner sep=1.1pt,
          font=\tiny\sffamily\bfseries, text=cAdd] {+};}%
}
\newcommand{\dop}[1][]{%
  \marginnote{\addedglyph\ifstrempty{#1}{}{\\[1pt]\scriptsize #1}}%
}

\newcommand{\legend}{%
  \begingroup\small\sffamily\color{inkSoft}
  \noindent{\color{ruleC}\rule{\textwidth}{0.5pt}}\\[4pt]
  \addedglyph\quad фрагмент дописан составителем: на лекции он не приводился
  или был дан устно.\par\vspace{2pt}
  \noindent{\color{ruleC}\rule{\textwidth}{0.5pt}}
  \endgroup\vspace{4pt}%
}

% ---------- русские имена для \cref --------------------------------
\makeatletter
\newcommand{\tcbcrefname}[3]{%
  \crefname{tcb@cnt@#1}{#2}{#3}\Crefname{tcb@cnt@#1}{\MakeUppercase#2}{\MakeUppercase#3}}
% счётчик один на все утверждения, поэтому \cref печатает только номер,
% а нужное слово в нужном падеже пишется руками: «по лемме~\ref{lem:...}»
\tcbcrefname{Theorem}{п.}{пп.}
\makeatother
\crefname{equation}{}{}
\crefname{figure}{рис.}{рис.}
\Crefname{figure}{Рис.}{Рис.}
\crefname{section}{§}{§§}
\Crefname{section}{§}{§§}
\crefname{algocf}{алгоритм}{алгоритмы}
\Crefname{algocf}{Алгоритм}{Алгоритмы}

% ---------- псевдокод ---------------------------------------------
\usepackage[ruled, vlined, linesnumbered, noend]{algorithm2e}
% Ключевые слова — стандартные английские (if / for / while / Input).
% Русифицирована только подпись; убрать — закомментировать строку ниже.
\SetAlgorithmName{Алгоритм}{алгоритм}{Список алгоритмов}
\SetAlFnt{\sffamily\small}
\SetAlCapFnt{\sffamily\bfseries\color{ink}}
\SetAlCapNameFnt{\sffamily\color{inkSoft}}
\SetAlgoNlRelativeSize{-1}
\SetNlSty{textnormal}{\color{inkSoft}}{}
\SetAlgoInsideSkip{smallskip}

% ---------- разное ------------------------------------------------
\usepackage{needspace}
\newcommand{\term}[1]{\textbf{#1}}          % вводимый термин
\newcommand{\eng}[1]{\textup{\textenglish{#1}}}

% ---------- шапка лекции ------------------------------------------
% \lecture{номер}{дата}{название}
\newcommand{\lecture}[3]{%
  \needspace{4\baselineskip}
  \section{#3}%
  \vspace{-3pt}
  {\footnotesize\sffamily\color{inkSoft}Лекция~#1 \quad$\cdot$\quad #2\par}
  \vspace{4pt}%
}

% Компактная шапка, когда лекция компилируется отдельным файлом
\newcommand{\standalonehead}{%
  \thispagestyle{plain}
  \begingroup\sffamily
  {\footnotesize\color{inkSoft}\subjectshort\ \ $\cdot$\ \ \lecturer\par}
  \vspace{2pt}
  {\color{ruleC}\rule{\textwidth}{0.8pt}}
  \endgroup\vspace{-8pt}%
}
